\scanpage{5}
$X$ \unsure{préschéma} [noethérien] \deleted{avec}\inserted{admettant}
Module inversible ample, $\dim X\leq1$

$K^\bullet(X)$ augmenté sur $H^0(X,Z)=Z^{\pi_0(X)}$, $E\longmapsto(\operatorname{rg}_i(E))_{i\in\pi_0(X)}$
\[
\left.\begin{gathered}
\text{L'idéal d'augmentation }\inserted{I(X)}\text{ est de carré nul}\\
\text{Si }x\in I(X),\text{ on a }\gamma^i(x)=0\text{ pour }i\geq2
\end{gathered}\right\}\quad F^i_\gamma(K^\bullet(X))=0\text{ si }i\geq2
\]
Donc
\[
K^\bullet(X)\simeq Z^{\pi_0(X)}\times K^1(X)
\]
\inserted{comme $\lambda$-anneau augmenté,}
\margin{$K^{(1)}(X)\simeq\operatorname{Pic}(X)=\prod_{i\in\pi_0(X)}\operatorname{Pic}(X_i)$
\inserted{$\simeq\deleted{\unsure{K}}\prod_{i\in\pi_0(X)}K^{(1)}(X_i)$}}

\deleted{\unsure{donc E}} $K^\bullet(X)$ entièrement déterminé (à isom près) par le système des
$\operatorname{Pic}(X_i)$, $i\in\pi_0(X)$

Si $E$ est un \deleted{fibré vectoriel}\inserted{faisceau loc libre} de rang $\unsure{n}$ sur $X$,
\unsure{on a donc|on a alors}
\[
\operatorname{cl}(E)=[\operatorname{rg}(E),c^1(E)]\qquad\text{où }c^1(E)=\det(E)
\]
\inserted{$\operatorname{rg}(E)=(\operatorname{rg}_i(E))_{i\in I}$}
\[
\left[\text{i.e.}\quad
\operatorname{cl}(E)=\sum_i\operatorname{rg}_i(E).1_i+[\operatorname{cl}(\det E)-1]
=\sum_i(\operatorname{rg}_i(E)-1)1_i+\operatorname{cl}(\det E)\right]
\]

\unsure{Car|Or}
\[
B^\bullet(X)=Z^{\pi_0(X)}\times\operatorname{Pic}(X)
\]
\diagram{\begin{tikzcd}
K^\bullet(X) \arrow[r,"c"] & B^\bullet(X)
\end{tikzcd}}
est l'identité, avec les identifications précédentes

\margin{$c$}

$K_\bullet(X)$ augmenté sur $Z^{\pi'_0(X)}$
\[
F\longmapsto(\operatorname{lg}_i(F))_{i\in\pi'_0(X)}
\]
\margin{$[\unsure{\pi^{(1)}(X)|\pi_0^{(1)}(X)}=$ ens des composantes irréductibles \uuline{de dim 1}$]$}

Le noyau est $K_0(X)$, il est isomorphe à
\[
Z^{\unsure{\pi_0^{(0)}(X)}}\times\underbrace{\operatorname{Pic}'(X)}
=Z^{\unsure{\pi_0^{(0)}(X)}}\times
\prod_{i\in\underbrace{\pi_0^{(1)}(X)}_{\substack{\text{ens des}\\\text{comp. conn.}\\\text{de }\uline{\text{dim 1}}}}}K_0(X_i)
\]

\margin{4}
